\clearpage
\VerdierPhysicalPageAnchor{0168}%
\begingroup\small\setlength{\parskip}{0pt}\setlength{\emergencystretch}{2.5em}
\VerdierSemanticAnchor{ch4-4-2-4-sequence}
\noindent\textit{The composite of the morphisms}
\[
E\!\begin{bmatrix}n-1\\ f_5,f_6,f_7\end{bmatrix}
\xrightarrow{d_1}
E\!\begin{bmatrix}n\\ f_3,f_4,f_5\end{bmatrix}
\xrightarrow{d_2}
E\!\begin{bmatrix}n+1\\ f_1,f_2,f_3\end{bmatrix}
\]
\textit{is zero. There is a canonical isomorphism}
\VerdierSemanticAnchor{ch4-4-2-4-1}
\[
\operatorname{Ker}(d_2)\big/\operatorname{Im}(d_1)
\simeq
E\!\begin{bmatrix}n\\ f_3f_2,f_4,f_6f_5\end{bmatrix}.
\tag{4.2.4.1}
\]
This follows immediately from (4.2.3.3) and (4.2.3.1).

\VerdierSemanticAnchor{ch4-4-2-5}
\medskip\noindent\textbf{4.2.5.}
The morphisms $d^n(\omega_5)$,
$\omega_5\in\operatorname{Ob}\mathcal{F}\ell_5(\mathcal J)$, are called
the \textit{differentials} (relative to the spectral object under
consideration). The family of functors (4.2.3.1)
\[
\omega_3\longmapsto
E\!\begin{bmatrix}n\\ \omega_3\end{bmatrix},
\qquad
\omega_3\in\operatorname{Ob}\mathcal{F}\ell_3(\mathcal J),
\]
equipped with the differentials (4.2.3.2) $d^n(\omega_5)$,
$\omega_5\in\operatorname{Ob}\mathcal{F}\ell_5(\mathcal J)$, is called
the \textit{spectral family} associated with the spectral object
$(H^n(f),\delta)$.

\VerdierSemanticAnchor{ch4-4-2-6}
\medskip\noindent\textbf{4.2.6.}
Let
$\omega_4=(\mathord\bullet\xrightarrow{f_1}\mathord\bullet
\xrightarrow{f_2}\mathord\bullet\xrightarrow{f_3}\mathord\bullet
\xrightarrow{f_4}\mathord\bullet)$
be an object of $\mathcal{F}\ell_4(\mathcal J)$. From it we obtain a
diagram in $\mathcal{F}\ell_3(\mathcal J)$:
\VerdierSemanticAnchor{ch4-4-2-6-fl3-diagram}
\[
\begin{tikzcd}[column sep=2.6em,row sep=2.15em]
{}\arrow[r,"\mathrm{id}"]\arrow[d,"f_1"']&
{}\arrow[r,"\mathrm{id}"]\arrow[d,"f_1"']&
{}\arrow[r,"\mathrm{id}"]\arrow[d,"f_1"']&
{}\arrow[r,"f_1"]\arrow[d,"f_2f_1"']&
{}\arrow[d,"f_2"]\\
{}\arrow[r,"\mathrm{id}"]\arrow[d,"f_2"']&
{}\arrow[r,"\mathrm{id}"]\arrow[d,"f_2"']&
{}\arrow[r,"f_2"]\arrow[d,"f_3f_2"']&
{}\arrow[r,"\mathrm{id}"]\arrow[d,"f_3"']&
{}\arrow[d,"f_3"]\\
{}\arrow[r,"\mathrm{id}"]\arrow[d,"f_3"']&
{}\arrow[r,"f_3"]\arrow[d,"f_4f_3"']&
{}\arrow[r,"\mathrm{id}"]\arrow[d,"f_4"']&
{}\arrow[r,"\mathrm{id}"]\arrow[d,"f_4"']&
{}\arrow[d,"f_4"]\\
{}\arrow[r,"f_4"]&
{}\arrow[r,"\mathrm{id}"]&
{}\arrow[r,"\mathrm{id}"]&
{}\arrow[r,"\mathrm{id}"]&{}
\end{tikzcd}
\]
Hence, on applying the functor
$E\!\begin{bmatrix}n\\ \omega_3\end{bmatrix}$, we obtain three exact
sequences in $\mathcal A$:
\VerdierSemanticAnchor{ch4-4-2-6-exact-sequences}
\[
0\longrightarrow
B\!\begin{bmatrix}n\\ f_2,f_4f_3\end{bmatrix}
 \Big/
B\!\begin{bmatrix}n\\ f_2,f_3\end{bmatrix}
\longrightarrow
E\!\begin{bmatrix}n\\ f_1,f_2,f_3\end{bmatrix}
\longrightarrow
E\!\begin{bmatrix}n\\ f_1,f_2,f_4f_3\end{bmatrix}
\longrightarrow0,
\]
\[
0\longrightarrow
E\!\begin{bmatrix}n\\ f_1,f_2,f_4f_3\end{bmatrix}
\longrightarrow
E\!\begin{bmatrix}n\\ f_1,f_3f_2,f_4\end{bmatrix}
\longrightarrow
E\!\begin{bmatrix}n\\ f_2f_1,f_3,f_4\end{bmatrix}
\longrightarrow0,
\]
\endgroup
